\documentclass[10pt,a4paper]{amsart}
\usepackage[margin=19mm]{geometry}
\usepackage{amsmath,amssymb,mathtools}
\usepackage[scr=rsfs,cal=euler]{mathalfa}
\usepackage{xfrac}
\usepackage[unicode,hidelinks,bookmarks=false]{hyperref}
\newcommand{\ie}{i.e.\ }
\newcommand{\cf}{cf.\ }
\newcommand{\resp}{respectively\ }
\newcommand{\Subsection}{\subsection}
\providecommand{\nobreakdash}{\nobreak\hbox{-}}
\setlength{\emergencystretch}{2em}
\multlinegap0pt
\usepackage{amssymb}
\usepackage{enumerate}
\usepackage{dsfont}
\usepackage{mathtools}
\usepackage{bbm}
\usepackage{tikz}
\usepackage{float}
\newtheorem{claim}{Claim}
\newtheorem{subcase}{Subcase}
\newtheorem{subsubcase}{Subsubcase}
\title{Off-diagonal Rado number for $x+y+c=z$ and $x+qy=z$}
\author{Rajat Adak, Yash Bakshi, L. Sunil Chandran, Saraswati Girish Nanoti}
\date{Source: arXiv:2602.23954v1 (CC BY 4.0)}
\begin{document}
\maketitle

\noindent\emph{Source and licence.} Off-diagonal Rado number for $x+y+c=z$ and $x+qy=z$. Version arXiv:2602.23954v1 (CC BY 4.0), distributed by arXiv under the Creative Commons Attribution 4.0 International licence, http://creativecommons.org/licenses/by/4.0/, as recorded in the arXiv licence metadata for this exact version and shown on the abstract page https://arxiv.org/abs/2602.23954v1. Full licence notice: \texttt{licenses/arxiv-2602.23954-CC-BY-4.0.txt}.\par\smallskip

\noindent\textit{Editorial scope.} Counted unit: the article's claim clm3.4 (source0.tex lines 326--328). The article's complete original proof of it is delivered with it (source0.tex lines 329--543, one \texttt{proof} environment of 215 lines). No mathematics is written, repaired or re-derived: the statement and the proof are the article's own lines, in the article's own notation, and no retained line is substituted or re-flowed.  No further source-local context is needed: every cross-reference of every retained line resolves inside the two delivered spans.  Nothing was omitted from any retained span, so the package carries no omission record.  No retained line contains a citation, so the wrapper carries no bibliography and no bibliography entry.  No retained line contains a figure environment, an image-inclusion command or drawing code, so no image is delivered and the package declares no asset dependency.  Editorial formatting only, in addition: an AMS \texttt{amsart} preamble that restates the article's own declarations and macros used by the retained lines, namely the environments \texttt{claim}, \texttt{subcase}, \texttt{subsubcase} on separate counters, together with the standard packages those lines need.  The retained environments print in this document as Claim 1, Subcase 1, Subsubcase 1 in order of appearance, whereas the article prints the same items with the numbering of its own full document.  The complete native source of the article as distributed in the arXiv e-print for this exact version is bundled unchanged as \texttt{sources/195436/source0.tex}.
\begin{claim}\label{clm3.4}
    $R_2(c,q) \leq (q+1)(c+2) + c +1$
\end{claim}

\begin{proof} Assume for the sake of contradiction, there is a coloring of $\{1,2,\dots,(q+1)(c+2)+c+1\}$ such that no red solution to $x+y+c=z$ and no blue solution to $x+y=z$ exists, respectively. We derive a
contradiction by analyzing the color assigned to $1$. Recall that, under the current case, we have either $c$ or $q$ to be even and $q >1$.

\begin{subcase}
    $c$ is even.
\end{subcase}
\begin{subsubcase} $1$ is colored red.\end{subsubcase}
\begingroup
\setlength{\jot}{2pt} % tighter vertical spacing inside align
\begin{align*}
(1,1,c+2)_c&\Longrightarrow {c+2} \text{ is blue}.\\
(c+2,c+2,(q+1)(c+2))_q&\Longrightarrow {(q+1)(c+2)}\text{ is red}.\\
((q+1)(c+2),1,(q+1)(c+2)+c+1)_c
  &\Longrightarrow {(q+1)(c+2)+c+1}\text{ is blue}.\\
(2c+3,c+2,(q+1)(c+2)+c+1)_q&\Longrightarrow {2c+3}\text{ is red}.\\
\Bigl(\tfrac{q(c+2)}{2}+1,\tfrac{q(c+2)}{2}+1,(q+1)(c+2)\Bigr)_c
  &\Longrightarrow {\tfrac{q(c+2)}{2}+1} \text{ is blue}.\\
(q(c+2)+1,1,(q+1)(c+2))_c&\Longrightarrow {q(c+2)+1} \text{ is blue}.\\
\Bigl(\tfrac{q(c+2)}{2}+1,\tfrac{c+2}{2},q(c+2)+1\Bigr)_q
  &\Longrightarrow {\tfrac{c+2}{2}}\text{ is red}.\\
\Bigl(\tfrac{c+2}{2},\tfrac{c+2}{2},2c+2\Bigr)_c&\Longrightarrow {2c+2} \text{ is blue}.\\
(2c+2,c+2,q(c+2)+2c+2)_q&\Longrightarrow {q(c+2)+2c+2} \text{ is red}. \displaybreak[2]\\
(q(c+2)+c+1,1,q(c+2)+2c+2)_c&\Longrightarrow {q(c+2)+c+1}\text{ is blue}.\\
(c+1,c+2,q(c+2)+c+1)_q&\Longrightarrow {c+1}\text{ is red} .\\
(c+1,2,2c+3)_c&\Longrightarrow {2}\text{ is blue}.\\
(2,2,2q+2)_q&\Longrightarrow {2q+2}\text{ is red}.\\
(2c+2,2,2q+2+2c)_q&\Longrightarrow {2q+2+2c}\text{ is red}.\\
(2q+2,c,2q+2+2c)_c&\Longrightarrow {c}\text{ is blue}.\\
(c,c,qc+c)_q&\Longrightarrow {qc+c}\text{ is red}.
\end{align*}
\text{Thus }$(2q+2,\ qc+c,\ q(c+2)+2c+2)$\text{ is a red solution, a contradiction.}
\endgroup

\begin{subsubcase}
    $1$ is colored blue.
\end{subsubcase}
\begingroup
\setlength{\jot}{2pt}
\begin{align*}
(1,1,q+1)_q&\Longrightarrow q+1 \text{ is red.}\\
(q+1,q+1,2q+2+c)_c&\Longrightarrow 2q+2+c \text{ is blue.}\\
(q+c+2,1,2q+2+c)_q&\Longrightarrow q+c+2 \text{ is red.}\\
(q+c+2,q+1,2q+2c+3)_c&\Longrightarrow 2q+2c+3 \text{ is blue.}\\
(2q+2c+3,1,3q+2c+3)_q&\Longrightarrow 3q+2c+3 \text{ is red.}\\
(2q+2+c,1,3q+c+2)_q&\Longrightarrow 3q+c+2 \text{ is red.}\\
(q+1,2q+1,3q+c+2)_c&\Longrightarrow 2q+1 \text{ is blue.}\\
(1,2,2q+1)_q&\Longrightarrow 2 \text{ is red.}\\
(2,2,c+4)_c&\Longrightarrow c+4 \text{ is blue}\\
(c+4,1,c+q+4)_q&\Longrightarrow c+q+4 \text{ is red.}\\
(q+1,3,c+q+4)_c&\Longrightarrow 3 \text{ is blue.}\\
(3,3,3q+3)_q&\Longrightarrow 3q+3 \text{ is red.}\\
(3q+3,c,3q+2c+3)_c&\Longrightarrow c \text{ is blue.}\\
(c,c,qc+c)_q&\Longrightarrow qc+c \text{ is red.}\\
(\tfrac{qc}{2},\tfrac{qc}{2},cq+c)_c&\Longrightarrow \tfrac{qc}{2} \text{ is blue.}\\
(\tfrac{qc}{2},1,\tfrac{qc}{2}+q)_q&\Longrightarrow \tfrac{qc}{2}+q \text{ is red.}\\
(\tfrac{qc}{2}+q,\tfrac{qc}{2}+q,qc+2q+c)_c&\Longrightarrow qc+2q+c \text{ is blue.}\\
(2q+c,c,qc+2q+c)_q&\Longrightarrow 2q+c \text{ is red.}\\
(3,1,q+3)_q&\Longrightarrow q+3 \text{ is red.}
\end{align*}
\text{Thus }$(q+3,\ 2q+c,\ 3q+2c+3)$\text{ is a red solution, a contradiction.}
\endgroup

\begin{subcase}
    $c$ is odd, and therefore $q$ must be even.
\end{subcase}


% \textbf{Case 2.1.} Assume $1$ is red.

% \begingroup
% \setlength{\jot}{2pt}
% \begin{align}
% &(1,1,c+2)\longrightarrow c+2 \text{ is blue.}\nonumber\\
% &(c+2,c+2,(q+1)(c+2))\longrightarrow (q+1)(c+2) \text{ is red.}\nonumber\\
% &\Bigl((q+1)(c+2),1,(q+1)(c+2)+c+1\Bigr)
% \longrightarrow (q+1)(c+2)+c+1 \text{ is blue.}\nonumber\\
% &(2c+3,c+2,(q+1)(c+2)+c+1)\longrightarrow 2c+3 \text{ is red.}\nonumber\\
% &\left(\tfrac{c+3}{2},\tfrac{c+3}{2},2c+3\right)
% \longrightarrow \tfrac{c+3}{2} \text{ is blue.}\nonumber\\
% &(qc+2q+1,1,(q+1)(c+2))\longrightarrow qc+2q+1 \text{ is blue.}\nonumber\\
% &\left(\tfrac{qc}{2}+q+1,\tfrac{qc}{2}+q+1,(q+1)(c+2)\right)
% \longrightarrow \tfrac{qc}{2}+q+1 \text{ is blue.}\nonumber\\
% &\left(\tfrac{qc}{2}+\tfrac{q}{2}+1,\tfrac{c+3}{2},qc+2q+1\right)
% \longrightarrow \tfrac{qc}{2}+\tfrac{q}{2}+1 \text{ is red.}\nonumber\\
% &\left(\tfrac{qc}{2}+\tfrac{q}{2}+1,1,\tfrac{qc}{2}+\tfrac{q}{2}+c+2\right)
% \longrightarrow \tfrac{qc}{2}+\tfrac{q}{2}+c+2 \text{ is blue.}\nonumber\\
% &(c+2,\tfrac{c+1}{2},\tfrac{qc}{2}+\tfrac{q}{2}+c+2)
% \longrightarrow \tfrac{c+1}{2} \text{ is red.}\nonumber\\
% &(1,\tfrac{c+1}{2},\tfrac{3(c+1)}{2})
% \longrightarrow \tfrac{3(c+1)}{2} \text{ is blue.}\nonumber\\
% &\left(\tfrac{c+5}{2},\tfrac{c+1}{2},2c+3\right)
% \longrightarrow \tfrac{c+5}{2} \text{ is blue.}\nonumber\\
% &\left(\tfrac{c+1}{2},\tfrac{c+1}{2},2c+1\right)
% \longrightarrow 2c+1 \text{ is blue.}\nonumber\\
% &\left(\tfrac{c+5}{2},\tfrac{c+5}{2},\tfrac{q(c+5)}{2}+\tfrac{c+5}{2}\right)
% \longrightarrow \tfrac{q(c+5)}{2}+\tfrac{c+5}{2} \text{ is red.}\nonumber\\
% &\left(\tfrac{3(c+1)}{2},\tfrac{c+5}{2},\tfrac{q(c+5)}{2}+\tfrac{3(c+1)}{2}\right)
% \longrightarrow \tfrac{q(c+5)}{2}+\tfrac{3(c+1)}{2} \text{ is red.}\nonumber\\
% &(2c+1,\tfrac{c+5}{2},\tfrac{q(c+5)}{2}+2c+1)
% \longrightarrow \tfrac{q(c+5)}{2}+2c+1 \text{ is red.}\nonumber\\
% &(1,\tfrac{q(c+5)}{2}+c,\tfrac{q(c+5)}{2}+2c+1)
% \longrightarrow \tfrac{q(c+5)}{2}+c \text{ is blue.}\nonumber\\
% &(c,\tfrac{c+5}{2},\tfrac{q(c+5)}{2}+c)
% \longrightarrow c \text{ is red.}\nonumber\\
% &(3,c,2c+3)\longrightarrow 3 \text{ is blue.}\nonumber\\
% &(3,\tfrac{c+5}{2},\tfrac{q(c+5)}{2}+3)
% \longrightarrow \tfrac{q(c+5)}{2}+3 \text{ is red.}\nonumber\\
% &(1,\tfrac{c+5}{2},\tfrac{q(c+5)}{2}+c+4)
% \longrightarrow \tfrac{q(c+5)}{2}+c+4 \text{ is blue.}\nonumber\\
% &(c+4,\tfrac{c+5}{2},\tfrac{q(c+5)}{2}+c+4)
% \longrightarrow c+4 \text{ is red.}\nonumber\\
% &(2,2,c+4)\longrightarrow 2 \text{ is blue.}\nonumber\\
% &(2,2,2q+4)\longrightarrow 2q+4 \text{ is red.}\nonumber\\
% &(2,\tfrac{c+5}{2},\tfrac{q(c+5)}{2}+2)
% \longrightarrow \tfrac{q(c+5)}{2}+2 \text{ is red.}\nonumber\\
% &(1,\tfrac{q(c+5)}{2}+2,\tfrac{q(c+5)}{2}+c+3)
% \longrightarrow \tfrac{q(c+5)}{2}+c+3 \text{ is blue.}\nonumber\\
% &(c+3,\tfrac{c+5}{2},\tfrac{q(c+5)}{2}+c+3)
% \longrightarrow c+3 \text{ is red.}\nonumber\\
% &(1,c+3,2c+4)\longrightarrow 2c+4 \text{ is blue.}\nonumber\\
% &(2c+4,2,2q+2c+4)\longrightarrow 2q+2c+4 \text{ is red.}\nonumber\\
% &(c+4,2q,2q+2c+4)\longrightarrow 2q \text{ is blue.}\nonumber\\
% &(2q,2,4q)\longrightarrow 4q \text{ is red.}\nonumber\\
% &(4q,1,4q+c+1)\longrightarrow 4q+c+1 \text{ is blue.}\nonumber\\
% &(q+c+1,3,4q+c+1)\longrightarrow q+c+1 \text{ is red.}\nonumber\\
% &(1,q,q+c+1)\longrightarrow q \text{ is blue.}\nonumber\\
% &(q,2,3q)\longrightarrow 3q \text{ is red.}\nonumber\\
% &(3q,1,3q+c+1)\longrightarrow 3q+c+1 \text{ is blue.}\nonumber\\
% &(c+1,3,3q+c+1)\longrightarrow c+1 \text{ is red.}\nonumber\\
% &(2q+c+1,2,4q+c+1)\longrightarrow 2q+c+1 \text{ is red.}\nonumber\\
% &(1,2q+c+1,2q+2c+2)\longrightarrow 2q+2c+2 \text{ is blue.}\nonumber\\
% &(2c+2,2,2q+2c+2)\longrightarrow 2c+2 \text{ is red.}
% \end{align}
% \endgroup

% \noindent
% Thus $(1,c+1,2c+2)$ is a red solution, a contradiction.

\begin{subsubcase}
    $1$ is colored red.
\end{subsubcase}

\begin{align*}
(1,1,c+2)_c&\Longrightarrow c+2 \text{ is blue.}\\
(c+2,c+2,(q+1)(c+2))_q&\Longrightarrow (q+1)(c+2) \text{ is red.}\\
((q+1)(c+2),1, (q+1)(c+2)+c+1 )_c&\Longrightarrow (q+1)(c+2)+c+1 \text{ is blue.}\\
(2c+3,c+2,(q+1)(c+2)+c+1)_q&\Longrightarrow 2c+3 \text{ is red.}\\
\left(\tfrac{c+3}{2},\tfrac{c+3}{2},2c+3\right)_c&\Longrightarrow \tfrac{c+3}{2} \text{ is blue.}\\
(qc+2q+1,1,(q+1)(c+2))_c&\Longrightarrow qc+2q+1 \text{ is blue.}\\
\left(\tfrac{qc}{2}+q+1,\tfrac{qc}{2}+q+1,(q+1)(c+2)\right)_c&\Longrightarrow \tfrac{qc}{2}+q+1 \text{ is blue.}\\
(\tfrac{qc}{2}+\tfrac{q}{2}+1,\tfrac{c+3}{2},qc+2q+1)_q&\Longrightarrow \tfrac{qc}{2}+\tfrac{q}{2}+1 \text{ is red.}\\
(\tfrac{qc}{2}+\tfrac{q}{2}+1,1,\tfrac{qc}{2}+\tfrac{q}{2}+c+2)_c&\Longrightarrow \tfrac{qc}{2}+\tfrac{q}{2}+c+2 \text{ is blue.}\\
(c+2,\tfrac{c+1}{2},\tfrac{qc}{2}+\tfrac{q}{2}+c+2)_q&\Longrightarrow \tfrac{c+1}{2} \text{ is red.}\\
(1,\tfrac{c+1}{2},\tfrac{3(c+1)}{2})_c&\Longrightarrow \tfrac{3(c+1)}{2} \text{ is blue.}\\
(\tfrac{c+1}{2}+2,\tfrac{c+1}{2},2c+3)_c&\Longrightarrow \tfrac{c+1}{2}+2=\tfrac{c+5}{2} \text{ is blue.}\\
(\tfrac{c+1}{2},\tfrac{c+1}{2},2c+1)_c&\Longrightarrow 2c+1 \text{ is blue.}\\
(\tfrac{c+5}{2},\tfrac{c+5}{2},\tfrac{q(c+5)}{2}+\tfrac{c+5}{2})_q&\Longrightarrow \tfrac{q(c+5)}{2}+\tfrac{c+5}{2} \text{ is red.}\\
(\tfrac{3(c+1)}{2},\tfrac{c+5}{2},\tfrac{q(c+5)}{2}+\tfrac{3(c+1)}{2})_q&\Longrightarrow \tfrac{q(c+5)}{2}+\tfrac{3(c+1)}{2} \text{ is red.}\\
(2c+1,\tfrac{c+5}{2},\tfrac{q(c+5)}{2}+2c+1)_q&\Longrightarrow \tfrac{q(c+5)}{2}+2c+1 \text{ is red.}\\
(1,\tfrac{q(c+5)}{2}+c,\tfrac{q(c+5)}{2}+2c+1)_c&\Longrightarrow \tfrac{q(c+5)}{2}+c \text{ is blue.}\\
(c,\tfrac{(c+5)}{2},\tfrac{q(c+5)}{2}+c)_q&\Longrightarrow c \text{ is red.}\\
(3,c,2c+3)_c&\Longrightarrow 3 \text{ is blue.}\\
(3,\tfrac{(c+5)}{2},\tfrac{q(c+5)}{2}+3)_q&\Longrightarrow \tfrac{q(c+5)}{2}+3 \text{ is red.}\\
(1,\tfrac{(c+5)}{2},\tfrac{q(c+5)}{2}+c+4)_c&\Longrightarrow \tfrac{q(c+5)}{2}+c+4\text{ is blue.}\\
(c+4,\tfrac{(c+5)}{2},\tfrac{q(c+5)}{2}+c+4)_q&\Longrightarrow c+4\text{ is red.}\\
(2,2,c+4)_c&\Longrightarrow 2 \text{ is blue.}\\
(2,2,2q+2)_q&\Longrightarrow 2q+2 \text{ is red.}\\
(2,\tfrac{c+5}{2},\tfrac{q(c+5)}{2}+2)_q&\Longrightarrow \tfrac{q(c+5)}{2}+2 \text{ is red.}\\
(1,\tfrac{q(c+5)}{2}+2,\tfrac{q(c+5)}{2}+c+3)_c&\Longrightarrow \tfrac{q(c+5)}{2}+c+3 \text{ is blue.}\\
(c+3,\tfrac{(c+5)}{2},\tfrac{q(c+5)}{2}+c+3)_q&\Longrightarrow c+3 \text{ is red.}\\
(1,c+3,2c+4)_c&\Longrightarrow 2c+4 \text{ is blue.}\\
(2c+4,2,2c+4)_q&\Longrightarrow 2q+2c+4 \text{ is red.}\\
(c+4,2q,2q+2c+4)_c&\Longrightarrow 2q \text{ is blue.}\\
(2q,2,4q)_q&\Longrightarrow 4q \text{ is red.}\\
(4q,1,4q+c+1)_c&\Longrightarrow 4q+c+1 \text{ is blue.}\\
(q+c+1,3,4q+c+1)_q&\Longrightarrow q+c+1 \text{ is red.}\\
(1,q,q+c+1)_c&\Longrightarrow q \text{ is blue.}\\
(q,2,3q)_q&\Longrightarrow 3q \text{ is red.}\\
(3q,1,3q+c+1)_c&\Longrightarrow 3q+c+1 \text{ is blue.}\\
(c+1,3,3q+c+1)_q&\Longrightarrow c+1 \text{ is red.}\\
(2q+c+1,2,4q+c+1)_q&\Longrightarrow 2q+c+1 \text{ is red.}\\
(1,2q+c+1,2q+2c+2)_c&\Longrightarrow 2q+2c+2 \text{ is blue.}\\
(2c+2,2,2q+2c+2)_q&\Longrightarrow 2c+2 \text{ is red.}
\end{align*}
Thus, $(1,c+1,2c+2)$ is a red solution, a contradiction.

\begin{subsubcase}
    $1$ is colored blue.
\end{subsubcase}
\begingroup
\setlength{\jot}{2pt}
\begin{align*}
(1,1,q+1)_q&\Longrightarrow q+1 \text{ is red.}\\
(q+1,q+1,2q+2+c)_c&\Longrightarrow 2q+2+c \text{ is blue.}\\
(q+c+2,1,2q+2+c)_q&\Longrightarrow q+c+2 \text{ is red.}\\
(q+c+2,q+1,2q+2c+3)_c&\Longrightarrow 2q+2c+3 \text{ is blue.}\\
(2q+2c+3,1,3q+2c+3)_q&\Longrightarrow 3q+2c+3 \text{ is red.}\\
(2q+2+c,1,3q+c+2)_q&\Longrightarrow 3q+c+2 \text{ is red.}\\
(q+1,2q+1,3q+c+2)_c&\Longrightarrow 2q+1 \text{ is blue.}\\
(1,2,2q+1)_q&\Longrightarrow 2 \text{ is red.}\\
(2,2,c+4)_c&\Longrightarrow c+4 \text{ is blue}\\
(c+4,1,c+q+4)_q&\Longrightarrow c+q+4 \text{ is red.}\\
(q+1,3,c+q+4)_c&\Longrightarrow 3 \text{ is blue.}\\
(3,3,3q+3)_q&\Longrightarrow 3q+3 \text{ is red.}\\
(3q+3,c,3q+2c+3)_c&\Longrightarrow c \text{ is blue.}\\
(c,c,qc+c)_q&\Longrightarrow qc+c \text{ is red.}\\
(\tfrac{qc}{2},\tfrac{qc}{2},cq+c)_c&\Longrightarrow \tfrac{qc}{2} \text{ is blue.}\\
(\tfrac{qc}{2},1,\tfrac{qc}{2}+q)_q&\Longrightarrow \tfrac{qc}{2}+q \text{ is red.}\\
(\tfrac{qc}{2}+q,\tfrac{qc}{2}+q,qc+2q+c)_c&\Longrightarrow qc+2q+c \text{ is blue.}\\
(2q+c,c,qc+2q+c)_q&\Longrightarrow 2q+c \text{ is red.}\\
(3,1,q+3)_q&\Longrightarrow q+3 \text{ is red.}
\end{align*}
\text{Thus }$(q+3,\ 2q+c,\ 3q+2c+3)$\text{ is a red solution, a contradiction.}
\endgroup
Thus, we arrive at a contradiction in both the subcases. Therefore, $R_2(c,q) \leq (q+1)(c+2)+c+1$. \end{proof}

\end{document}
